\documentclass[12pt,a4paper]{amsart}

\usepackage{amsmath,amssymb,amsthm}
\usepackage{mathtools}
\usepackage{enumitem}
\usepackage{hyperref}

\newtheorem{theorem}{Theorem}[section]
\newtheorem{lemma}[theorem]{Lemma}
\newtheorem{proposition}[theorem]{Proposition}
\newtheorem{corollary}[theorem]{Corollary}
\theoremstyle{definition}
\newtheorem{definition}[theorem]{Definition}
\newtheorem{assumption}[theorem]{Assumption}
\newtheorem{problem}[theorem]{Open Problem}
\theoremstyle{remark}
\newtheorem{remark}[theorem]{Remark}

\newcommand{\Rig}{\textup{[Rig]}}
\newcommand{\Cond}{\textup{[Cond]}}
\newcommand{\Open}{\textup{[Open]}}

\begin{document}

\title[Lagrangian Singularity Transport under Fourier Duality]{%
  Lagrangian Singularity Transport under Fourier Duality:\\
  Cubic WKB Phases, A$_2$-Fold Singularities,\\
  and the Airy Normal Form}

\author{[Author]}

\begin{abstract}
This paper works in the \emph{one-dimensional setting}
($\beta,u\in\mathbb{R}$).
A cubic WKB phase $c_3\beta^3$ produces an A$_2$-fold
singularity of the canonical relation (Arnol'd \cite{Arnold1972}).
Setting $h := c^{-1} \to 0$, we work semiclassically throughout;
all estimates are uniform in $h$ in the sense of rescaled
phase space $(\beta,\xi)\sim(1,h^{-1})$.

The analysis rests on two theorems and one proposition
forming an acyclic proof chain:
Theorem~\ref{thm:CFU-geom} (geometric): reduces an A$_2$-fold
phase to Airy normal form via canonical transformation;
Theorem~\ref{thm:CFU-symbol} (analytic): transports the conormal
amplitude $|\beta|^{-1/2}\in I^{-1/4,0}(\Lambda_0,\Lambda_1)$
over the already-reduced Airy phase to
$a_{\mathrm{CFU}}\in S^{-1/2}_{1,0}(h)$ uniformly in $h$;
Proposition~\ref{prop:airy-id} (identification): identifies the
model integral with the Airy function microlocally.

Three consequences:
(I)~$|\Phi(u;h)|\le C(1+|u|)^{-3/4}$ uniformly in $h$;
(II)~$|\int_0^U\Phi|=O(U^{1/4})$;
(III)~A$_2$-stability for PSWF as $h\to 0$ remains open
(Open Problem~\ref{prob:PSWF}).
\end{abstract}

\maketitle

\section{Definitions}
\label{sec:defs}

\begin{definition}[A$_2$-fold singularity]\label{def:A2}
Let $\mathcal{C}\subset T^*\mathbb{R}_\beta\times T^*\mathbb{R}_u$
be Lagrangian, $\pi_u:\mathcal{C}\to T^*\mathbb{R}_u$ the projection.
$\mathcal{C}$ has an \emph{A$_2$-fold at $p_0$} if
$\det d\pi_u|_{p_0}=0$ and $d(\det d\pi_u)|_{p_0}\ne 0$
(Arnol'd \cite{Arnold1972}).
For phase $\Psi(\beta,u)=\Theta(\beta;h^{-1})+u\beta$
this reduces to $\Theta''(0)=0$, $\Theta'''(0)\ne 0$.
\end{definition}

\begin{definition}[Conormal amplitude class]\label{def:Iml}
The Melrose--Uhlmann class
$I^{m,l}(\Lambda_0,\Lambda_1)$ \cite{MelroseU1979}
consists of distributions with conormal singularities
at intersecting Lagrangians.
The amplitude $|\beta|^{-1/2}\chi$ belongs to
$I^{-1/4,0}(\Lambda_0,\Lambda_1)$,
where $\Lambda_0$ is the zero-section and
$\Lambda_1$ the conormal to $\{\beta=0\}$.
Standard FIO class $I^m$ is insufficient here
since the amplitude singularity coincides with
the critical point.
\end{definition}

\begin{definition}[Semiclassical parameter]\label{def:h}
Set $h:=c^{-1}$, $h\to 0$ as $c\to\infty$.
All symbol estimates below are uniform in $h\in(0,h_0]$
in the sense of rescaled phase space
$(\beta,\xi)\sim(1,h^{-1})$, i.e.\
$h$-dependent rescaling is implicit in the
$S^m_{1,0}(h)$-norm:
\[
  S^m_{1,0}(h) := \bigl\{a(u;h)\in C^\infty(\mathbb{R}_u)
  \;\big|\;
  |\partial_u^\alpha a|\le C_\alpha(1+|u|)^{m-|\alpha|},\;
  \text{uniformly in }h\in(0,h_0]\bigr\}
\]
(Zworski \cite{Zworski2012}, Ch.~4).
\end{definition}

\section{Two CFU Theorems and the Airy Identification}
\label{sec:CFU}

\begin{assumption}[A$_2$-fold conditions, \Cond]
\label{ass:A2}
\begin{enumerate}[\rm(A1)]
\item \textit{(Single nondegenerate fold.)}
  $\Theta(\beta;h^{-1})=c_3(h)\beta^3+O(\beta^4)$
  with $c_3(h)\ge\delta_0>0$ uniformly in $h$,
  and $\Theta''(\beta)\ne 0$ for $\beta\ne 0$ near $0$.
\item \textit{(CFU admissibility.)}
  The CFU change of variables $\beta\mapsto t(\beta,u;h)$
  has $|\det(dt/d\beta)|\ge\delta_1>0$
  with $|\partial^\alpha_{\beta,u}t|\le C_\alpha$
  uniformly in $h$.
\item \textit{(Unique fold.)}
  $\beta=0$ is the only critical point on
  $\mathrm{supp}\,\chi$.
\end{enumerate}
\end{assumption}

\begin{theorem}[CFU: geometric normal form, \Rig]
\label{thm:CFU-geom}
Under (A1) and (A3),
there exists a smooth local canonical transformation
$\kappa:(\beta,\xi)\mapsto(t,\tau)$
in a fixed conic neighbourhood of the fold point $p_0$,
a real function $\mu(u;h)$ with $|\mu|\sim|u|^{1/2}$
uniformly in $h$,
and a real phase $R(u;h)$,
such that
\[
  \Theta(\beta;h^{-1})+u\beta
  = \tfrac{1}{3}t^3-\mu(u;h)\,t+R(u;h)
  +O(t^\infty).
\]
This is the \emph{Airy phase normal form} of the A$_2$-fold.
This theorem involves only the phase function;
no amplitude or integral is referenced here.
\end{theorem}

\begin{proof}
Standard CFU reduction \cite{CFU1957}:
(A1) gives $\Theta''(0)=0$, $\Theta'''(0)\ne 0$,
so the phase is locally equivalent to
$\frac{1}{3}t^3-\mu t$ by Morse-type normal form theory
(Golubitsky--Guillemin \cite{Golub1973}).
The scaling $|\mu|\sim|u|^{1/2}$ follows from the
fold geometry.
Uniformity in $h$ follows from $c_3(h)\ge\delta_0>0$ (A1).
\end{proof}

\begin{theorem}[CFU: symbol transport, \Rig]
\label{thm:CFU-symbol}
Under Assumption~\ref{ass:A2},
with the phase already in Airy normal form
(Theorem~\ref{thm:CFU-geom}),
the amplitude $|\beta|^{-1/2}\chi
\in I^{-1/4,0}(\Lambda_0,\Lambda_1)$
(Definition~\ref{def:Iml})
is transported by the canonical transformation $\kappa$
as follows:
\begin{enumerate}[\rm(i)]
\item The $I^{m,l}$-class is preserved under $\kappa$
  (Melrose--Uhlmann \cite{MelroseU1979};
  Guillemin--Uhlmann \cite{GuillemU1981}, Lemma~2.1).
\item Restriction of the transported distribution to
  the stationary set $t=0$ yields the
  \emph{transported principal symbol}
  $a_{\mathrm{CFU}}(u;h)\in S^{-1/2}_{1,0}(h)$,
  satisfying $|a_{\mathrm{CFU}}(u;h)|\le C|u|^{-1/2}$
  uniformly in $h$.
\item The fold-zone integral satisfies,
  modulo $S^{-\infty}$ uniformly in $h$:
  \begin{equation}\label{eq:CFU}
    \Phi_{\mathrm{fold}}(u;h)
    = e^{i(\alpha_*u+R(u;h))}
    \bigl[
      a_{\mathrm{CFU}}(u;h)\cdot F(u;h)
      + r(u;h)
    \bigr],
  \end{equation}
  where $F(u;h)$ is the Airy model integral
  (Proposition~\ref{prop:airy-id})
  and $|r(u;h)|\le C_N(1+|u|)^{-N}$ for all $N$.
\end{enumerate}
Schematically:
$\mathrm{CFU}:(\mathcal{C},I^{-1/4,0})
\longrightarrow
(\text{Airy phase},S^{-1/2}(h),S^{-\infty}).$
\end{theorem}

\begin{proof}
(i): (A2) + \cite{MelroseU1979,GuillemU1981}.
(ii): The $I^{-1/4,0}$-order maps to $S^{-1/2}$
under restriction to the stationary set,
by the principal symbol calculus for $I^{m,l}$-distributions.
(iii): Combining (i)--(ii) with Theorem~\ref{thm:CFU-geom}
and absorbing the Maslov factor into $R$
(H\"ormander \cite{Hormander1983}, Ch.~25.1).
Remainder bounds from the CFU asymptotic expansion \cite{CFU1957}.
\end{proof}

\begin{proposition}[Airy model identification, \Rig]
\label{prop:airy-id}
The Airy model integral
\[
  F(u;h) :=
  \int_\mathbb{R}
  e^{i(\frac{1}{3}t^3-\mu(u;h)\,t)}
  \,dt
\]
satisfies $F(u;h)=\mathrm{Ai}(\mu(u;h))$
(up to a universal constant factor).
For $|\mu|\ge 1$,
$|\mathrm{Ai}(\mu)|\le C|\mu|^{-1/4}$
(Olver \cite{Olver1974}, Ch.~11).
This identification is purely a special function fact;
it is the only point where the Airy function enters
the argument.
\end{proposition}

\section{Main Results}
\label{sec:results}

\begin{lemma}[Degenerate stationary phase
  with conormal weight, \Rig]
\label{lem:degenerate-phase}
For $I(\mu):=\int_\mathbb{R}|t|^{-1/2}
e^{i(\frac{1}{3}t^3-\mu t)}\,dt$,
$|\mu|\ge 1$:
$|I(\mu)|\le C|\mu|^{-3/4}$.
\end{lemma}

\begin{proof}
By Theorems~\ref{thm:CFU-geom}--\ref{thm:CFU-symbol}
and Proposition~\ref{prop:airy-id},
\[
  I(\mu) = a_{\mathrm{CFU}}(\mu)\cdot\mathrm{Ai}(\mu) + r(\mu)
\]
with $|a_{\mathrm{CFU}}|\le C'|\mu|^{-1/2}$
(Theorem~\ref{thm:CFU-symbol}(ii))
and $|r|\le C_N|\mu|^{-N}$.
By Proposition~\ref{prop:airy-id},
$|\mathrm{Ai}(\mu)|\le C|\mu|^{-1/4}$.
Hence
\[
  |I(\mu)|
  \le C'|\mu|^{-1/2}\cdot C|\mu|^{-1/4}
  + C_N|\mu|^{-N}
  = C''|\mu|^{-3/4}.
\]
The exponent $-3/4=(-1/2)+(-1/4)$:
$(-1/2)$ from conormal symbol order
(Theorem~\ref{thm:CFU-symbol}(ii));
$(-1/4)$ from A$_2$-fold Airy decay
(Proposition~\ref{prop:airy-id}).
Both are determined by the same fold geometry.
\end{proof}

\begin{corollary}[Pointwise decay, \Cond]
\label{cor:pointwise}
Under Assumption~\ref{ass:A2},
$|\Phi(u;h)|\le C(1+|u|)^{-3/4}$
uniformly in $h\in(0,h_0]$.
\end{corollary}

\begin{proof}
Fold zone: Theorems~\ref{thm:CFU-geom},
\ref{thm:CFU-symbol},
Proposition~\ref{prop:airy-id},
Lemma~\ref{lem:degenerate-phase}.
Far zone ($|\partial_\beta^2\Psi|\ge c_0'>0$):
standard stationary phase gives $O(|u|^{-1/2})$.
Compact zone $|u|\le 1$: continuity.
\end{proof}

\begin{corollary}[Ces\`aro improvement, \Cond]
\label{cor:cesaro}
Under Assumption~\ref{ass:A2},
$|\int_0^U\Phi(u;h)\,du|=O(U^{1/4})$,
improving $O(U^{1/2})$ by pure $L^1$-integrability
of $(1+|u|)^{-3/4}$.
\end{corollary}

\subsection{Converse and rigidity}

\begin{assumption}[Inverse reconstruction,
  conjectural, \Cond]
\label{hyp:rigidity}
\textit{Forward (rigorous):}
Theorems~\ref{thm:CFU-geom}--\ref{thm:CFU-symbol}
give A$_2$-fold $\Rightarrow$ Airy class.
\textit{Inverse (conjectural):}
If $\Phi$ is in the Airy class,
$\sigma(g\mapsto\Phi)|_\mathcal{C}\ne 0$,
$g\mapsto\Phi$ injective mod smoothing,
and $\mathrm{WF}(\Phi)\subseteq\pi_u(\mathcal{C})$
on the A$_2$ stratum,
then $\mathcal{C}$ has an A$_2$-fold uniformly in $h$.
Verification for PSWF: Paper~XVII.
\end{assumption}

\begin{corollary}[Conditional equivalence]
\label{cor:equiv}
Under Assumption~\ref{hyp:rigidity}:
Airy normal form $\Rightarrow$ Assumption~\ref{ass:A2}.
The forward direction holds without
Assumption~\ref{hyp:rigidity}.
\end{corollary}

\section{Open Problem}
\label{sec:open}

\begin{problem}[A$_2$-stability for PSWF, \Open]
\label{prob:PSWF}
Does the PSWF amplitude satisfy
Assumption~\ref{ass:A2} uniformly as $h\to 0$?
(a)~$c_3(h)\ge\delta_0$: expected from Olver \cite{Olver1974};
uniform remainder bounds required.
(b)~Unique fold uniformly: plausible, not established.
(c)~CFU admissibility (A2) uniformly in $h$:
technically hardest; no result currently available.
\end{problem}

\begin{thebibliography}{9}
\bibitem{PaperXV}
[Author], \textit{Paper~XV}, 2026.
\bibitem{Olver1974}
F.~W.~J.~Olver,
\textit{Asymptotics and Special Functions},
Academic Press, 1974.
\bibitem{Zworski2012}
M.~Zworski,
\textit{Semiclassical Analysis}, AMS, 2012.
\bibitem{Hormander1983}
L.~H\"ormander,
\textit{The Analysis of Linear PDE~I},
Springer, 1983.
\bibitem{CFU1957}
C.~Chester, B.~Friedman, F.~Ursell,
Proc.\ Cambridge Philos.\ Soc.\
\textbf{53} (1957), 599--611.
\bibitem{Arnold1972}
V.~I.~Arnol'd,
Funct.\ Anal.\ Appl.\
\textbf{6} (1972), 254--272.
\bibitem{Golub1973}
M.~Golubitsky, V.~Guillemin,
\textit{Stable Mappings}, Springer, 1973.
\bibitem{GuillemU1981}
V.~Guillemin, G.~Uhlmann,
Duke Math.\ J.\ \textbf{48} (1981), 251--267.
\bibitem{MelroseU1979}
R.~Melrose, G.~Uhlmann,
Comm.\ Pure Appl.\ Math.\
\textbf{32} (1979), 483--519.
\end{thebibliography}

\end{document}
